%! Principal Component Analysis — Unit 7, Lesson 7.3 — Student Packet Cover Sheet
\documentclass[10pt]{article}
\usepackage{linalg-article}
\usepackage{linalg-boxes}
\usepackage{ltablex}
\keepXColumns

\begin{document}

% Full-bleed burgundy banner
\begin{tikzpicture}[remember picture, overlay]
  \fill[burgundy] (current page.north west)
    rectangle ([yshift=-0.9in]current page.north east);
\end{tikzpicture}
\vspace{-0.8in}

\begin{center}
  {\color{white}\textbf{\LARGE Linear Algebra}}\\[4pt]
  {\color{blushmid}\large\textbf{Unit 7 \quad The Singular Value Decomposition}}\\[6pt]
  {\color{white}\textbf{\large Lesson 7.3 \quad Principal Component Analysis}}
\end{center}

\vspace{0.22in}
\namedateperiod
\vspace{0.18in}

\begin{learningtargetbox}
\small
\begin{enumerate}[leftmargin=1.5em, itemsep=5pt]
  \item \ldots{} \textbf{center} a data table and form $A^{\top}A$, then read its eigenvectors as the \textbf{principal components} $\mathbf{v}_1\perp\mathbf{v}_2$ (the directions the data spreads along, biggest first).
  \item \ldots{} measure the fraction of the spread the first component captures ($\sigma_1^2/(\sigma_1^2+\sigma_2^2)$) and interpret what PC1 \emph{means} for the data.
  \item \ldots{} explain why keeping the top few components summarizes data in far fewer numbers --- the SVD's ``keep the biggest $\sigma$'' idea applied to data.
\end{enumerate}
\end{learningtargetbox}

\vspace{0.14in}

\begin{tocbox}
\small
\renewcommand{\arraystretch}{1.5}
\begin{tabularx}{\linewidth}{@{} c l X r @{}}
\rowcolor{blush}
\textbf{\#} & \textbf{Component} & \textbf{Description} & \textbf{Score} \\
\midrule
1 & Warm-Up        & Spiral review: center a data table, form $A^{\top}A$, and find its eigenvalues       & \blank{1.2cm} \\
2 & Guided Notes   & Centering, principal components as eigenvectors of $A^{\top}A$, and reading the spread & \blank{1.2cm} \\
3 & Group Activity & Full PCA on a data table, the spread fraction, and choosing how many components to keep & \blank{1.2cm} \\
4 & Exit Ticket    & Quick check: find PC1, the spread fraction, and justify                              & \blank{1.2cm} \\
5 & Homework       & PCA end to end, spread percentages, a line-shaped data set, and a spread spectrum      & \blank{1.2cm} \\
\midrule
  & \multicolumn{2}{r}{\textbf{Total}} & \blank{1.2cm} \\
\end{tabularx}
\end{tocbox}

\vspace{0.10in}

\begin{remindbox}
\small \textbf{Principal Component Analysis} finds the directions data spreads along the most. Write the data
as a matrix, \textbf{center} it (subtract each column's mean), and form the symmetric matrix $A^{\top}A$. Its
eigenvectors are the \textbf{principal components} $\mathbf{v}_1\perp\mathbf{v}_2$ --- perpendicular
directions ordered by how much the data spreads, biggest first --- and its eigenvalues are $\sigma^2$, the
spread along each. The first component PC1 is the best single direction through the data cloud; keeping the
top few components summarizes the data in far fewer numbers, just as keeping the biggest singular values
compressed an image in Lesson~7.2.
\end{remindbox}

\end{document}
